Summanden müssen dieselbe Einheit haben, und beide Seiten einer Gleichung ebenfalls.
\begin{abcliste}
\abc $p^2$ hat die Einheit von $v^2$, also $[p] = \result{\si{\meter\per\second}}$. Aus $[a]\cdot[q] = \si{\meter\squared\per\second\squared}$ folgt $[q] = \result{\si{\meter}}$.
\abc Umgeformt ist $p^2 = \dfrac{2s}{a}$ mit der Einheit $\dfrac{\si{\meter}}{\si{\meter\per\second\squared}} = \si{\second\squared}$, also $[p] = \result{\si{\second}}$.
\abc $p$ wird von $a$ subtrahiert: $[p] = \result{\si{\meter\per\second\squared}}$. Damit $[s] = \si{\meter}$ ist, muss $[q] = \result{\si{\second\squared}}$ sein.
\end{abcliste}
